\section{Corrección de la transcripción y normalización de términos}

Este episodio no tiene subtítulos de la plataforma; el texto se basa en una transcripción hecha con faster-whisper large-v3 en CUDA. Los términos de matemáticas y de IA son abundantes, y la transcripción automática produce errores sistemáticos: por ejemplo, escribe Lean como «令» o «林», escribe AI for Math con palabras de sonido parecido, oye mal nombres propios como Carina, Hong o Verve Coffee y fragmenta términos como theorem proving, proof assistant y formalization. Al generar los apuntes es obligatorio corregir los términos; de lo contrario, el lector malinterpretará la línea técnica central.

\subsection{Términos clave: tabla de corrección de errores comunes de transcripción}

\noindent\begin{longtable}{p{0.24\textwidth}p{0.32\textwidth}p{0.35\textwidth}}
\toprule
Término normalizado & Posibles errores de transcripción & Motivo de la corrección \\
\midrule
Carina Hong / Hong Letong (洪乐潼) & 洪乐同, 洪乐童 y grafías erróneas de Carina por sonido parecido & El nombre está ligado a los antecedentes de Axiom; debe unificarse. \\
Axiom & 公理 («axioma»), A型 («tipo A»), Axium & Nombre de la empresa; en chino puede explicarse como «axioma», pero en el texto debe conservarse en inglés. \\
Lean & 令, 林, line & Nombre del proof assistant; está ligado a la línea central de la matemática formalizada. \\
AI for Math & Al for Math, AI for Mark & Dirección central de este episodio. \\
Proof Assistant & 证明助手 («asistente de demostración») y proof system usados indistintamente & Se refiere a sistemas de verificación como Lean o Coq. \\
Formalization & Confusión entre 形式化 («formalización») y 格式化 («dar formato») & Se refiere a escribir las matemáticas en un sistema formal, no a dar formato tipográfico. \\
Ken Ono & 小野肯, 小野 Ken & Nombre propio de un matemático, relacionado con las señales sobre la organización de Axiom. \\
Bounded / Free Attention & Traducciones erróneas como «límite de atención» o «atención libre» & Aquí es una metáfora de modos de pensamiento matemático, no el mecanismo de attention del Transformer en sí. \\
\bottomrule
\end{longtable}

\begin{warningbox}{La transcripción automática no puede usarse directamente como texto técnico}
El material de entrevistas es especialmente propenso a errores de transcripción en nombres propios, títulos de artículos y nombres de sistemas. Unas notas de alto estándar deben corregir los términos a partir de los capítulos, del contexto y de fuentes públicas externas; no pueden copiar palabra por palabra los subtítulos automáticos.
\end{warningbox}

\subsection{Resumen de la sección}

La transcripción es solo materia prima, no son los apuntes. Un contenido con tantos términos como AI for Math requiere especialmente normalizar primero los términos y después escribir notas estructuradas.
